\providecommand{\permil}{\mathord{\text{‰}}}
\paragrafos{Ποσοστά: το κλάσμα με παρονομαστή το εκατό}{Α.5.1, «Ποσοστά»}

\begin{anagke}
Στο σχολείο γίνεται ψηφοφορία για την επιτροπή της εκδρομής και η Φανή είναι υποψήφια. Τα τρία τμήματα της πρώτης τάξης ψηφίζουν χωριστά και τα αποτελέσματα, γραμμένα στον πίνακα, είναι: στο Α1 η Φανή πήρε $12$ ψήφους από τους $20$ που ψήφισαν, στο Α2 πήρε $17$ από τους $25$ και στο Α3 πήρε $21$ από τους $35$.

«Στο Α3 με ψήφισαν περισσότεροι», λέει η Φανή περήφανα, «είκοσι ένας!». Ο Νίκος σηκώνει το φρύδι: «Στο Α3 όμως ψήφισαν πολύ περισσότερα παιδιά. Δεν συγκρίνεις το ίδιο πράγμα». Ο Βασίλης προτείνει να τα συγκρίνουν σαν κλάσματα, και η Πολυξένη ρωτάει πώς θα το έλεγαν πιο γρήγορα, με μία ματιά: «Στις ειδήσεις λένε πάντα "το 60\%" ή "το 12\%". Μήπως μπορούμε να κάνουμε όλα τα κλάσματα να έχουν παρονομαστή το εκατό, για να συγκρίνονται εύκολα;».

Το ίδιο ισχύει στη ζωή. Ο Ηλίας είδε στη βιτρίνα «έκπτωση 30\%» και η θεία Μαρία διάβασε σε μια απόδειξη «ΦΠΑ 24\%». Αν είναι να μιλάμε για μέρη ενός συνόλου, χρειαζόμαστε μια κοινή μονάδα σύγκρισης, και η πιο συνηθισμένη είναι το «τοις εκατό».
\end{anagke}

\begin{lexi}{ποσοστό, τοις εκατό}{ποσοστό, τοις εκατό}
Η λέξη \lexiword{ποσοστό} έρχεται από το «πόσο»: ένα \emph{πόσο} μέρος του συνόλου. Το «\lexiword{τοις εκατό}» σημαίνει «ανά εκατό»: το $12\%$ σημαίνει «$12$ για κάθε $100$». Το σύμβολο $\%$ θυμίζει ένα κλάσμα, με το $1$ και το $0$ και $0$ να έχουν κατέβει στη θέση του παρονομαστή. Το \lexiword{τοις χιλίοις} ($\permil$) σημαίνει «ανά χίλια».
\end{lexi}

\begin{metafrasi}
Δώδεκα για κάθε εκατό. & $12\%=\dfrac{12}{100}$\\
Το $35\%$ των $80$ ευρώ. & $\dfrac{35}{100}\cdot80=28$\\
Τρία από τα πέντε ίσα μέρη. & $\dfrac{3}{5}=\dfrac{60}{100}=60\%$\\
Το ποσοστό ως δεκαδικός. & $37\%=0{,}37$\\
Τέσσερα για κάθε χίλια. & $4\permil=\dfrac{4}{1\,000}$\\
\end{metafrasi}

\tmima{Η θεωρία}

\begin{orismos}{ποσοστό επί τοις εκατό}
Το \textbf{ποσοστό $\alpha\%$} (διαβάζεται «$\alpha$ τοις εκατό») είναι ο αριθμός $\dfrac{\alpha}{100}$. Το \textbf{$\alpha\permil$} («$\alpha$ τοις χιλίοις») είναι ο αριθμός $\dfrac{\alpha}{1\,000}$. Το \textbf{$\alpha\%$ του $\beta$} είναι ο αριθμός $\dfrac{\alpha}{100}\cdot\beta$.
\end{orismos}

Άρα κάθε ποσοστό είναι ένα κλάσμα (ή δεκαδικός), και οι κανόνες των κλασμάτων και των δεκαδικών ισχύουν. Μόνο που τώρα, αντί να λέμε «τα $\frac{3}{4}$ των $80$», λέμε «το $75\%$ των $80$».

\tmimasmall{Από κλάσμα σε ποσοστό και ανάποδα}

\begin{itemize}[leftmargin=1.2em,itemsep=0.3ex,topsep=2pt]
\item \textbf{Κλάσμα σε ποσοστό:} κάνω το κλάσμα να έχει παρονομαστή $100$ (ή το διαιρώ και βρίσκω δεκαδικό), και μετά γράφω το αριθμητή με το σύμβολο $\%$. Για παράδειγμα $\frac{4}{5}=\frac{80}{100}=80\%$ και $\frac{3}{8}=0{,}375=37{,}5\%$.
\item \textbf{Δεκαδικός σε ποσοστό:} πολλαπλασιάζω με το $100$ (η υποδιαστολή πάει δύο θέσεις δεξιά): $0{,}52=52\%$, $3{,}41=341\%$, $0{,}07=7\%$.
\item \textbf{Ποσοστό σε κλάσμα:} γράφω $\dfrac{\alpha}{100}$ και απλοποιώ: $12\%=\frac{12}{100}=\frac{3}{25}$.
\item \textbf{Ποιο ποσοστό είναι το $\mu$ του $\nu$:} διαιρώ $\mu:\nu$ και μετατρέπω σε ποσοστό: $15$ € από τα $60$ € είναι $\frac{15}{60}=0{,}25=25\%$.
\end{itemize}

Ένα ποσοστό μπορεί να είναι και μεγαλύτερο από $100\%$: το $150\%$ ενός αριθμού είναι ο ίδιος αριθμός μαζί με το μισό του. Και το $100\%$ ενός αριθμού είναι όλος ο αριθμός.

\tmimasmall{Το ποσοστό στην καθημερινή ζωή}

Πολλές εκφράσεις της καθημερινότητας είναι ποσοστά: «Αύξηση $14\%$» (στον μισθό $100$ € προστίθενται $14$ €), «έκπτωση $60\%$» (από κάθε $100$ € δεν πληρώνουμε $60$), «ΦΠΑ $24\%$» (στην αξία $100$ € προστίθενται $24$ € φόρος), «επιτόκιο $2\%$», «είμαι $100\%$ σίγουρος» (χωρίς καμία αμφιβολία), «πιθανότητα $1\%$» (μια φορά στις εκατό). Για να βρούμε το ποσοστό μιας ποσότητας, την πολλαπλασιάζουμε με το κλάσμα του ποσοστού.

\begin{prosoxi}
\begin{itemize}[leftmargin=1.2em,itemsep=0.3ex,topsep=2pt]
\item Το $\%$ δεν είναι αριθμός από μόνο του: το $30\%$ είναι πάντα το $30\%$ \emph{κάποιου} συνόλου. Το $30\%$ των $50$ δεν είναι ίσο με το $30\%$ των $200$.
\item Όταν συγκρίνω ποσοστά, πρέπει να είναι του ίδιου συνόλου ή να μιλάω για ποσοστά ως τιμές: τα $21$ από τα $35$ δεν είναι περισσότερα «σαν ποσοστό» από τα $17$ από τα $25$, παρότι $21>17$.
\item Στη μετατροπή δεκαδικού σε ποσοστό μετακινώ την υποδιαστολή \emph{δύο} θέσεις δεξιά και όχι μία: $0{,}3=30\%$, όχι $3\%$.
\item Το $\permil$ είναι δέκα φορές μικρότερο από το $\%$: $1\%=10\permil$.
\end{itemize}
\end{prosoxi}

\tmima{Λυμένα παραδείγματα}

\paradeigma{Τα ποσοστά της Φανής}
\tixero Η Φανή πήρε $12$ από $20$ ψήφους στο Α1, $17$ από $25$ στο Α2 και $21$ από $35$ στο Α3.\par
\tithelo Σε ποιο τμήμα την ψήφισαν «περισσότερο».\par
\possindeo Βρίσκω το ποσοστό κάθε τμήματος: $\frac{12}{20}=\frac{60}{100}=60\%$, $\frac{17}{25}=\frac{68}{100}=68\%$ και $\frac{21}{35}=\frac{3}{5}=\frac{60}{100}=60\%$.\par
\apantisi Στο Α2, με $68\%$. Στα Α1 και Α3 πήρε από $60\%$.\par
\elegxos Στο Α3 πήρε περισσότερες ψήφους ($21$), αλλά ψήφισαν και περισσότεροι. Το ποσοστό δίνει το σωστό μέτρο.

\paradeigma{Μετατροπές}
\tixero Τα κλάσματα $\frac{4}{5}$, $\frac{3}{8}$, $\frac{5}{7}$ και τα ποσοστά $15\%$, $64\%$, $2{,}5\%$.\par
\tithelo Τα κλάσματα ως ποσοστά (το $\frac{5}{7}$ στο δέκατο του ποσοστού) και τα ποσοστά ως ανάγωγα κλάσματα.\par
\possindeo $\frac{4}{5}=\frac{80}{100}=80\%$. $\frac{3}{8}=0{,}375=37{,}5\%$. $\frac{5}{7}=0{,}7142\ldots\approx0{,}714=71{,}4\%$. Και $15\%=\frac{15}{100}=\frac{3}{20}$, $64\%=\frac{64}{100}=\frac{16}{25}$, $2{,}5\%=\frac{2{,}5}{100}=\frac{25}{1\,000}=\frac{1}{40}$.\par
\apantisi $80\%$, $37{,}5\%$, περίπου $71{,}4\%$ και $\frac{3}{20}$, $\frac{16}{25}$, $\frac{1}{40}$.\par
\elegxos $\frac{1}{40}=0{,}025=2{,}5\%$.

\paradeigma{Το ποσοστό ενός ποσού}
\tixero Η Ρήνη βλέπει ένα μπλουζάκι των $80$ € με έκπτωση $35\%$ και ένα βιβλίο των $50$ € με ΦΠΑ $24\%$.\par
\tithelo Το ποσό της έκπτωσης και η νέα τιμή, και ο ΦΠΑ και η τιμή με ΦΠΑ.\par
\possindeo Έκπτωση $35\%$ του $80$: $\frac{35}{100}\cdot80=28$ €, νέα τιμή $80-28=52$ €. ΦΠΑ $24\%$ του $50$: $\frac{24}{100}\cdot50=12$ €, τιμή $50+12=62$ €.\par
\apantisi $52$ € και $62$ €.\par
\elegxos Το $35\%$ είναι λίγο περισσότερο από το $\frac{1}{3}$, και το $\frac{1}{3}$ του $80$ είναι περίπου $27$. Συμφωνεί.

\paradeigma{Ποιο ποσοστό είναι;}
\tixero Ο Νίκος έβαλε $14$ από $20$ βολές, ο Βασίλης $9$ από $12$ και ο Στέλιος $5$ από $8$.\par
\tithelo Ποιος έχει το καλύτερο ποσοστό επιτυχίας.\par
\possindeo $\frac{14}{20}=0{,}7=70\%$, $\frac{9}{12}=0{,}75=75\%$, $\frac{5}{8}=0{,}625=62{,}5\%$.\par
\apantisi Ο Βασίλης με $75\%$.\par
\elegxos Ο Νίκος έβαλε τις περισσότερες βολές (14), αλλά ο Βασίλης έχει το καλύτερο ποσοστό.

\paradeigma{Το τοις χιλίοις}
\tixero Ένα εργοστάσιο βγάζει $25\,000$ τεμάχια και τα $4\permil$ είναι ελαττωματικά.\par
\tithelo Πόσα ελαττωματικά είναι.\par
\possindeo $4\permil=\frac{4}{1\,000}$. Άρα $\frac{4}{1\,000}\cdot25\,000=4\cdot25=100$.\par
\apantisi $100$ τεμάχια.\par
\elegxos $100$ από $25\,000$ είναι $\frac{100}{25\,000}=0{,}004=4\permil$.

\tmima{Ασκήσεις}

\zoni{Το κατάλαβα}{μετατροπές και πράξεις με ποσοστά}

\ex Γράψε ως ποσοστά επί τοις εκατό τα κλάσματα: α) $\frac{1}{5}$ \quad β) $\frac{3}{2}$ \quad γ) $\frac{1}{4}$ \quad δ) $\frac{3}{4}$ \quad ε) $\frac{7}{20}$.

\ex Γράψε ως ποσοστά τους δεκαδικούς: α) $0{,}37$ \quad β) $2{,}5$ \quad γ) $0{,}08$ \quad δ) $0{,}6$ \quad ε) $0{,}005$.

\ex Γράψε ως ανάγωγα κλάσματα τα ποσοστά: α) $15\%$ \quad β) $8\%$ \quad γ) $45\%$ \quad δ) $60\%$ \quad ε) $125\%$ \quad στ) $0{,}5\%$.

\ex Υπολόγισε: α) το $10\%$ των $3\,500$ € \quad β) το $25\%$ των $64$ κιλών \quad γ) το $40\%$ της ώρας \quad δ) το $75\%$ του λίτρου \quad ε) το $5\%$ των $240$ \quad στ) το $150\%$ των $30$.

\ex Τι ποσοστό είναι: α) τα $15$ € από τα $60$ € \quad β) τα $45$ λεπτά από την ώρα \quad γ) τα $7$ από τα $28$ \quad δ) τα $9$ παιδιά από τα $24$ \quad ε) τα $330$ από τα $1\,100$.

\ex α) Γράψε το $7\permil$ ως δεκαδικό και ως κλάσμα. β) Βρες το $25\permil$ του $4\,000$. γ) Πόσα τοις χιλίοις είναι το $1\%$;

\ex Σωστό ή λάθος; α) $50\%=\frac{1}{2}$ \quad β) Το $100\%$ ενός αριθμού είναι ο ίδιος ο αριθμός. \quad γ) Το $200\%$ του $30$ είναι $60$. \quad δ) Το $10\%$ του $50$ ισούται με το $50\%$ του $10$. \quad ε) $\frac{1}{4}=14\%$. \quad στ) $0{,}3\%=0{,}3$.

\zoni{Το χρησιμοποιώ}{ψηφοφορίες, μπάσκετ, εκπτώσεις και φόροι}

\ex \textbf{Η ψηφοφορία.} Στη δεύτερη ψηφοφορία ο Βασίλης πήρε $15$ από $25$ ψήφους στο Α1, $24$ από $40$ στο Α2 και $27$ από $36$ στο Α3. Βρες τα ποσοστά και πες σε ποιο τμήμα τον ψήφισαν περισσότερο.

\ex \textbf{Το μπάσκετ.} Η Αγγελική έβαλε $18$ από $24$ βολές και η Ρήνη $13$ από $20$. Ποια έχει το καλύτερο ποσοστό;

\ex \textbf{Τα αθλητικά.} Ένα ζευγάρι αθλητικά κοστίζει $80$ € και έχει έκπτωση $25\%$. Πόσο είναι η έκπτωση και πόσο η νέα τιμή;

\ex \textbf{Το ΦΠΑ.} Το ποδήλατο κοστίζει $150$ € χωρίς ΦΠΑ. Πόσο πληρώνει ο Ηλίας αν ο ΦΠΑ είναι $24\%$;

\ex \textbf{Το μπουκάλι.} Η Ζωή ήπιε το $40\%$ από ένα μπουκάλι $1{,}5$ λίτρων. Πόσο ήπιε και πόσο έμεινε;

\ex \textbf{Το κέρδος του κυλικείου.} Το κυλικείο έβγαλε $450$ € και το $20\%$ ήταν καθαρό κέρδος. Πόσο ήταν το κέρδος;

\ex \textbf{Η τάξη.} Στην Α3 υπάρχουν $25$ παιδιά και τα $11$ είναι αγόρια. Τι ποσοστό της τάξης είναι τα αγόρια και τι τα κορίτσια;

\ex \textbf{Το νερό.} Το νερό είναι περίπου το $60\%$ του βάρους ενός ανθρώπου. Πόσα κιλά νερό υπάρχουν σε ένα παιδί βάρους $45$ κιλών;

\zoni{Το σκέφτομαι}{ποσοστά που ξεγελάνε}

\ex Δώσε ένα παράδειγμα ποσοστού μεγαλύτερου από το $100\%$. Μια τιμή ανέβηκε από $50$ € σε $110$ €. Η νέα τιμή είναι τι ποσοστό της παλιάς;

\ex Ποιο είναι μεγαλύτερο, το $40\%$ του $50$ ή το $50\%$ του $40$; Εξήγησε γιατί.

\ex Μια τιμή των $100$ € αυξάνεται κατά $10\%$ και μετά μειώνεται κατά $10\%$. Επιστρέφει στα $100$ €;

\ex Πώς βρίσκεις γρήγορα με το μυαλό το $10\%$ και το $1\%$ ενός αριθμού; Βρες το $10\%$ και το $1\%$ του $3\,450$ και μετά το $11\%$.

\ex Ο Στέλιος λέει «είμαι $80\%$ σίγουρος ότι θα βρέξει». Τι εννοεί; Και τι θα σήμαινε «$100\%$ σίγουρος», «$0\%$ σίγουρος»;

\zoni{Για το σπίτι}{ποσοστά στα νέα}

\ex Βρες πέντε φράσεις με ποσοστά από εφημερίδα ή ειδήσεις και εξήγησε τι σημαίνει η καθεμία.

\ex Πάρε μια απόδειξη σούπερ μάρκετ και βρες το ποσό του ΦΠΑ που αναγράφεται. Τι ποσοστό της συνολικής αξίας είναι;
